% Generated draft: requires mathematical and source review.
\documentclass[11pt]{article}
\usepackage[margin=25mm]{geometry}
\usepackage{fontspec}
\setmainfont{Latin Modern Roman}
\usepackage{amsmath,amssymb,mathtools,unicode-math}
\setmathfont{Latin Modern Math}
\usepackage{enumitem,needspace,xurl,hyperref}
\setlength{\parindent}{0pt}
\setlength{\parskip}{5pt}
\setlength{\emergencystretch}{3em}
\newcommand{\N}{\mathbb{N}}
\newcommand{\Z}{\mathbb{Z}}
\newcommand{\Q}{\mathbb{Q}}
\newcommand{\R}{\mathbb{R}}
\newcommand{\C}{\mathbb{C}}
\newcommand{\F}{\mathbb{F}}
\newcommand{\A}{\mathbb{A}}
\newcommand{\PP}{\mathbb{P}}
\newcommand{\E}{\mathbb{E}}
\newcommand{\eps}{\varepsilon}
\newcommand{\dd}{\,\mathrm{d}}
\begin{document}
\section*{UnsolvedMath ID: 30002452}
\textbf{Source problem code:} catalogue entry
\subsection{Definitions and mathematical statement}
A Poisson--Lie group is a group equipped with a Poisson structure for which multiplication is a Poisson map, with the product Poisson structure on the Cartesian square. The supplied record does not specify the group \(G\), its ground field, or the category in which this definition is to be applied.

In the usual root-theoretic terminology, after fixing a set \(\Delta\) of simple roots and its invariant inner product, a Belavin--Drinfeld triple is a triple \(T=(\Gamma_1,\Gamma_2,\tau)\), where \(\Gamma_1,\Gamma_2\subseteq\Delta\), the bijection \(\tau:\Gamma_1\to\Gamma_2\) preserves inner products, and every \(\alpha\in\Gamma_1\) eventually leaves \(\Gamma_1\) under iteration of \(\tau\). More explicitly, for each such \(\alpha\) there is an integer \(m\geq1\) such that \(\tau^j(\alpha)\in\Gamma_1\) for \(0\leq j<m\), while \(\tau^m(\alpha)\notin\Gamma_1\). This recalls the conventional terminology; the root data and admissible class of groups intended by the source remain unspecified.

\textbf{Source question.} For every Belavin--Drinfeld triple \(T\) associated with the intended group \(G\), does its double \(D(G)\) admit a generalized cluster structure compatible with the corresponding Poisson--Lie structure?

Writing \(\operatorname{BD}(G)\) for the intended collection of triples and \(\pi_T\) for the intended corresponding Poisson structure on the double, the quantifier pattern is
\[
\forall T\in\operatorname{BD}(G)\quad
\exists\,\mathcal C_T\quad
\bigl[\mathcal C_T\text{ is a generalized cluster structure on }D(G)
\text{ compatible with }\pi_T\bigr].
\]
Here \(D(G)\), \(\pi_T\), and the required generalized cluster structures denote the objects intended by the supplied source, rather than objects completely determined by the record.

\textbf{Unresolved qualifications.} The record does not define the double, specify whether additional continuous parameters or choices accompany \(T\), or give the axioms and coefficient conventions for the generalized cluster structure. It also does not specify which algebra of functions on \(D(G)\) is to carry that structure, whether regularity or equality with an upper generalized cluster algebra is required, or the precise meaning of compatibility. One common coordinate interpretation of Poisson compatibility requires every extended cluster \((x_1,\ldots,x_N)\) to satisfy
\[
\{x_i,x_j\}_{\pi_T}=\omega_{ij}x_ix_j
\]
with coefficients \(\omega_{ij}\) constant on the underlying space; the supplied record does not establish whether this is the entire intended condition. Consequently, the displayed question preserves the source target but requires these conventions to be recovered before it becomes a fully specified existence assertion.
\subsection{Short English statement}
The question asks whether every Belavin--Drinfeld triple gives rise to a compatible generalized cluster structure on the intended Poisson--Lie double of \(G\). The supplied record leaves the group, double, Poisson parameters, and cluster conventions unspecified. Its label \texttt{partially\_solved} and its claims about known cases are supplied dated curation, not independently verified mathematical conclusions.
\subsection{Sources}
\begin{itemize}
\item Supplied local UnsolvedMath registry record, permanent ID \(30002452\), title: \emph{Generalized Cluster Structures on Poisson--Lie Doubles}; source problem code \texttt{catalogue entry}; collection: \emph{Oberwolfach Reports Open Problems}. No individual proposer is identified; the supplied proposed year is 2013.
\item Supplied source citation: \emph{Cluster Algebras and Related Topics} (2014), Oberwolfach Reports, pp.~3379--3432, DOI \texttt{10.4171/owr/2013/58}. Source URL and literature-source label ``original OWR report'': \url{https://doi.org/10.4171/owr/2013/58}.
\item Supplied literature-source label ``Zhang generalized double'': \url{https://doi.org/10.1080/00927872.2025.2450036}. The record supplies no fuller bibliographic citation or theorem reference.
\item Literature assessment dated exactly \texttt{2026-08-21}: the supplied curation labels the problem partially solved and reports constructions for trivial and Cremmer--Gervais data and later isolated groups, while reporting that existence for every Belavin--Drinfeld triple remains unproved. This is supplied dated curation, not independently verified; the record itself requests independent verification of time-sensitive status and cited claims.
\item The record labels the statement \texttt{corrected\_verified} and reports that it was checked against the official Oberwolfach report. This verification claim is supplied provenance; no independent source consultation is asserted here.
\end{itemize}
\end{document}
